\section{Vatican II: Event, Texts, and Contested Renewal}

The Second Vatican Council opened on 11 October 1962 with more than 2,400 fathers from a Church far less European than Vatican I. Orthodox and Protestant observers were present; journalists reported globally; episcopal conferences and transnational theological networks coordinated interventions. John XXIII did not supply a finished reform platform. The council acquired its direction through decisions about agenda, drafting, and language made by the fathers under papal authority.

\subsection{The first session changes the process}

The Curia's preparatory schemas expected rapid approval. The bishops instead delayed election of commission members so that they could consult, rejected the draft on Revelation after a major conflict, and returned other texts for reconstruction. The liturgy schema, already the fruit of decades of papally supervised work, advanced most readily. It supplied the council's first constitution and a practical test of reform.

This procedural redirection took place within papal authority. John XXIII permitted it; after his death on 3 June 1963, Paul VI reconvened the council, named its purposes, moderated its work, intervened in difficult texts, and promulgated every document. The episcopal college acted with and under its head through visible debate.

\subsection{Four sessions and the principal contests}

The second session (29 September--4 December 1963) promulgated \emph{Sacrosanctum Concilium} and \emph{Inter mirifica}. Debate on the Church established the sequence that became \emph{Lumen gentium}: mystery, People of God, hierarchy, laity, universal holiness, religious, eschatology, and Mary. The placement of “People of God” before hierarchy expressed baptismal priority without abolishing office.

The third session (14 September--21 November 1964) produced \emph{Lumen gentium}, \emph{Unitatis redintegratio}, and \emph{Orientalium Ecclesiarum}. Collegiality occasioned intense dispute. The appended Preliminary Note clarified that the college never exists without its head and does not derive authority from a democratic delegation. Paul VI also reserved or influenced certain questions, including the title used for Mary and the institutional form of the Synod of Bishops. “Black Week” is partisan shorthand for these late interventions; the promulgated text and note, not the emotional label, govern doctrine.

The fourth session (14 September--8 December 1965) completed eleven documents. \emph{Dei Verbum} resolved the Revelation schema through a Christ-centered account of divine self-communication, Scripture and Tradition flowing from one source, inspiration, truth, interpretation, and the Church's teaching office. \emph{Dignitatis humanae} passed after repeated revision and a last postponement; it grounded civil immunity from coercion in personal dignity while expressly leaving intact the duty toward truth and the traditional doctrine of the true religion. \emph{Nostra aetate} rejected collective Jewish guilt and antisemitism while treating non-Christian religions briefly. \emph{Gaudium et spes} addressed humanity in a deliberately time-conscious pastoral constitution.

\subsection{The sixteen documents}

Titles do not by themselves assign every sentence one theological grade. “Dogmatic constitution” identifies principal doctrinal subject and form; “pastoral constitution” can contain doctrine; a declaration can teach doctrine and state prudential applications.

\begin{longtable}{>{\bfseries\raggedright\arraybackslash}p{0.24\linewidth}>{\raggedright\arraybackslash}p{0.15\linewidth}>{\raggedright\arraybackslash}p{0.16\linewidth}>{\raggedright\arraybackslash}p{0.37\linewidth}}
\toprule
\textbf{Document} & \textbf{Genre} & \textbf{Promulgated} & \textbf{Principal object}\\
\midrule
\endhead
\emph{Sacrosanctum Concilium} & Constitution & 4 Dec 1963 & Sacred liturgy: nature, participation, reform principles, sacraments, Office, year, music, and art.\\
\emph{Inter mirifica} & Decree & 4 Dec 1963 & Social communications, moral responsibility, and ecclesial media.\\
\emph{Lumen gentium} & Dogmatic constitution & 21 Nov 1964 & Mystery and constitution of the Church, People of God, hierarchy, laity, holiness, religious, eschatology, and Mary.\\
\emph{Orientalium Ecclesiarum} & Decree & 21 Nov 1964 & Dignity, discipline, sacraments, and patrimony of Eastern Catholic Churches.\\
\emph{Unitatis redintegratio} & Decree & 21 Nov 1964 & Catholic principles and practice of ecumenism.\\
\emph{Christus Dominus} & Decree & 28 Oct 1965 & Bishops' pastoral office, dioceses, cooperation, and conferences.\\
\emph{Perfectae caritatis} & Decree & 28 Oct 1965 & Renewal and adaptation of religious life from Gospel, founder, sound tradition, and present need.\\
\emph{Optatam totius} & Decree & 28 Oct 1965 & Priestly formation.\\
\emph{Gravissimum educationis} & Declaration & 28 Oct 1965 & Christian education, parents, schools, and universities.\\
\emph{Nostra aetate} & Declaration & 28 Oct 1965 & Relations with non-Christian religions; rejection of collective Jewish guilt, antisemitism, and religious discrimination.\\
\emph{Dei Verbum} & Dogmatic constitution & 18 Nov 1965 & Divine Revelation, its transmission, inspiration, interpretation, and Scripture in Church life.\\
\emph{Apostolicam actuositatem} & Decree & 18 Nov 1965 & Vocation and forms of the lay apostolate.\\
\emph{Dignitatis humanae} & Declaration & 7 Dec 1965 & Civil religious freedom as immunity from coercion, with continuing duties toward truth.\\
\emph{Ad gentes} & Decree & 7 Dec 1965 & Missionary nature and activity of the Church.\\
\emph{Presbyterorum ordinis} & Decree & 7 Dec 1965 & Ministry and life of priests.\\
\emph{Gaudium et spes} & Pastoral constitution & 7 Dec 1965 & Human dignity, the modern world, marriage, culture, economics, politics, peace, and the Church's mutual relation with history.\\
\bottomrule
\end{longtable}

\subsection{A new style, not a new faith}

Earlier councils often used short canons and anathemas to exclude errors. Vatican II normally used extended biblical, patristic, sacramental, and pastoral exposition. John O'Malley describes a rhetoric of invitation, dialogue, participation, service, friendship, and pilgrimage. Style matters: it proposes a way of exercising authority and seeing the Church's relation to the world. It does not make the content merely literary.

The change also creates interpretive difficulty. A canon condemning one proposition can be mapped narrowly. A long pastoral paragraph may weave doctrine, exhortation, historical diagnosis, and program. The reader must ask what subject is taught, with what verbs, from which sources, to what end, and whether a claim repeats prior definitive teaching or applies it prudentially.

The council's great doctrinal syntheses include Revelation as God's self-gift fulfilled in Christ; the Church as sacrament, communion, People and Body; episcopal consecration and collegiality; the universal call to holiness; the genuine vocation of the laity; the real though imperfect communion with the Catholic Church of those who believe in Christ and have been truly baptized (\emph{Unitatis redintegratio} 3); the dignity of conscience without relativism; and the missionary orientation of the whole Church. Many elements were ancient or already developing. Their conciliar arrangement changed Catholic self-understanding and pastoral priority.

\subsection{Communism and Masonry: silence, analysis, and continuity}

Vatican II did not name either Communism or Freemasonry in a promulgated document. That fact requires explanation, not invention. Hundreds of fathers sought an explicit treatment of Communism, but the final \emph{Gaudium et spes} analyzed systematic atheism by its claims: humanity as its own sole end and maker, economic emancipation used against religion, state coercion, and the suppression of faith in education. Footnotes referenced prior condemnations. The council's chosen method was diagnosis and dialogue with persons, not silence about the ideology's evil.

The prior teachings and disciplinary judgments represented by \emph{Divini Redemptoris}, the Masonic prohibition, and the Creator--creature distinction remained in force. The council addressed systematic atheism through diagnosis and dialogue while leaving the existing Masonic discipline undisturbed. Its omission of the proper nouns nevertheless contributed to a perception that older conflicts had simply disappeared. The pastoral strategy can be judged by its documented choices, clarity, and fruits.

\subsection{Closing and promulgation}

Paul VI solemnly promulgated each text after the fathers' votes. On 7 December 1965 the mutual Catholic--Orthodox declarations lifted the memory of the 1054 excommunications from the midst of the Churches without restoring full communion. On 8 December, the apostolic brief \emph{In Spiritu Sancto} closed the council and ordered observance of its decisions.

The legal and ceremonial solemnity is undeniable. The remaining question is theological: did those solemn sessions make every conciliar sentence an infallible definition, no doctrine at all, or something differentiated? The council itself answered.

\judgmentbox{Vatican II was a genuine exercise of the Church's supreme magisterium whose fathers deliberately changed the inherited drafts, literary form, and pastoral strategy. Those changes were neither a new religion nor a null event. The promulgated texts, not an imagined “spirit” and not the preparatory schemas they rejected, are the council. Their authority must be assigned claim by claim.}
